% Clase 13 --- El circuito de dos mallas (§4.3).
% El circuito que el docente copió de su cuaderno. Lo que la figura tiene que
% dejar claro son las dos decisiones que se toman ANTES de escribir ecuaciones:
% inventar un sentido para cada corriente (si sale negativo, iba al revés) y
% elegir dos mallas independientes --- la tercera no aporta información nueva.
\begin{center}
\begin{tikzpicture}[scale=0.92]

  \begin{scope}[line width=0.9pt, color=figblue]
    % rama izquierda
    \draw (0,0) to[battery1, l=$\varepsilon_1$] (0,2.8);
    \draw (0,2.8) to[R, l=$R_1$] (3.4,2.8);
    % rama central
    \draw (3.4,2.8) to[R, l_=$R_3$] (3.4,0);
    % rama derecha
    \draw (3.4,2.8) to[R, l=$R_2$] (6.8,2.8);
    \draw (6.8,2.8) to[battery1, l=$\varepsilon_2$] (6.8,0);
    % rama inferior
    \draw (0,0) -- (6.8,0);
    \fill (3.4,2.8) circle (2.2pt);
    \fill (3.4,0) circle (2.2pt);
  \end{scope}

  % corrientes: sentidos inventados
  \draw[figvecres] (0.7,2.55) -- (1.7,2.55);
  \node[figlbl, figred, anchor=north] at (1.2,2.5) {$I_1$};
  \draw[figvecres] (6.1,2.55) -- (5.1,2.55);
  \node[figlbl, figred, anchor=north] at (5.6,2.5) {$I_2$};
  \draw[figvecres] (3.85,1.85) -- (3.85,0.95);
  \node[figlbl, figred, anchor=west] at (3.9,1.4) {$I_3$};

  % mallas independientes
  \draw[figamber, line width=0.9pt, dashed] (0.65,0.5) rectangle (2.85,2.0);
  \node[figlbl, anchor=north] at (1.75,1.9) {\footnotesize malla 1};
  \draw[figamber, line width=0.9pt, dashed] (4.15,0.5) rectangle (6.35,2.0);
  \node[figlbl, anchor=north] at (5.25,1.9) {\footnotesize malla 2};

  % nudo
  \node[figlblaux, anchor=north] at (3.4,-0.15) {nudo: $I_3 = I_1 + I_2$};

  % la malla exterior, que no aporta nada nuevo
  \draw[figgray, line width=0.6pt, dotted] (-0.5,-0.75) rectangle (7.3,3.45);
  \node[figlblaux, anchor=north, align=center] at (3.4,-0.95)
    {la malla exterior también es un camino cerrado, pero es la suma\\[-0.2em]
     de las otras dos: su ecuación no agrega nada};

\end{tikzpicture}\\[0.5em]
{\small\itshape Con más de una fuente ya no se puede saber de antemano hacia
dónde circula cada corriente, porque las dos baterías empujan en sentidos
opuestos y gana la que gana según los valores relativos. La salida no es
adivinar: se les asigna un sentido cualquiera, se resuelve, y si alguna sale
\textbf{negativa} es que circula al revés que la flecha dibujada ---la
convención se corrige sola---. La segunda decisión es cuántas ecuaciones de
malla escribir: las dos marcadas alcanzan, y con la de nudos queda un sistema
lineal $3\times3$ en $I_1$, $I_2$, $I_3$. Los signos al recorrer una malla: una
resistencia atravesada a favor de su corriente resta $R\,I$ y en contra suma;
una FEM recorrida en el sentido de su flecha suma $\varepsilon$ y en contra
resta.}
\end{center}
